% GENERATED FILE. Edit canonical lessons, then run npm run generate:books.
% canonical-source-hash: fnv1a64:4589df619659bea1
% canonical-lessons: TE-C448-disa, TE-C448-sucika, TE-C448-dari-tappu, TE-C448-malupu-tirugu, TE-C448-addadari

\chapter{Finding the Way}
\label{ch:te-finding-the-way}
\hlchaptermodality{448}

\begin{chapteropening}
\textbf{By the end of this chapter:} \emph{I can say a direction, a signpost, to lose my way, to take a turn and a shortcut.}

Complete the last lesson of chapter 448: I can say a direction, a signpost, to lose my way, to take a turn and a shortcut.
\end{chapteropening}

\section[\texorpdfstring{\te{దిశ} (diśa)}{దిశ (diśa)}]{\te{దిశ} (diśa) — a compass direction; a side}

\label{lesson:TE-C448-disa}



\begin{quote}
Before the new one: say the Telugu for a footpath, then the Telugu for to break a rule.
\end{quote}

\subsection*{You'll want to know: \te{దిశ}}
\textbf{\te{దిశ}} — \emph{diśa} — \textquotedblleft{}a compass direction; a side\textquotedblright{}.

\begin{quote}
\te{సూర్యుడు} \te{తూర్పు} \te{దిశలో} \te{ఉదయిస్తాడు}.
\end{quote}

\begin{quote}
\emph{sūryuḍu tūrpu diśalō udayistāḍu.}
\end{quote}

\begin{quote}
The sun rises in the east.
\end{quote}

The four directions are \te{తూర్పు}, \te{పడమర}, \te{ఉత్తరం} and \te{దక్షిణం}. \te{దిక్కు} is a more everyday word.

\begin{cousinweb}
\textbf{\te{దిశ}} (\emph{diśa}) is borrowed from Sanskrit \emph{diśā}, \textquotedblleft{}direction, region\textquotedblright{} (CDIAL 6339).
\end{cousinweb}

\begin{grammarlens}[title={What you've built}]
One more word for getting around.
\end{grammarlens}

\subsection*{Your turn}
\begin{itemize}
\raggedright
  \item \emph{Say it:} \emph{diśa}
  \item \emph{Say it:} \emph{diśa}, once more
  \item \emph{Say it:} \emph{sūryuḍu tūrpu diśalō udayistāḍu.} — The sun rises in the east.
  \item \emph{Recall:} say the Telugu for a footpath, then the Telugu for to break a rule, then say \emph{diśa} again
\end{itemize}

\begin{tcolorbox}[breakable,colback=teal!4,colframe=teal!35!black,title={Before you move on}]
What does \te{దిశ} mean? (\textquotedblleft{}A compass direction; a side\textquotedblright{}.) Say it once more.
\end{tcolorbox}
\section[\texorpdfstring{\te{సూచిక} (sūcika)}{సూచిక (sūcika)}]{\te{సూచిక} (sūcika) — a signpost; an indicator}

\label{lesson:TE-C448-sucika}



\begin{quote}
Before the new one: say the Telugu for to break a rule, then the Telugu for a compass direction.
\end{quote}

\subsection*{You'll want to know: \te{సూచిక}}
\textbf{\te{సూచిక}} — \emph{sūcika} — \textquotedblleft{}a signpost; an indicator\textquotedblright{}.

\begin{quote}
\te{ఆ} \te{సూచిక} \te{మీద} \te{నగరానికి} \te{దూరం} \te{రాసి} \te{ఉంది}.
\end{quote}

\begin{quote}
\emph{ā sūcika mīda nagarāniki dūraṁ rāsi undi.}
\end{quote}

\begin{quote}
The distance to the city is written on that signpost.
\end{quote}

Any sign that tells you something: a signpost, an index in a book, a gauge on a dashboard.

\begin{cousinweb}
\textbf{\te{సూచిక}} (\emph{sūcika}) is built from Sanskrit \emph{sūc}, \textquotedblleft{}to indicate, to make known\textquotedblright{}, whose noun \emph{sūcaka} means \textquotedblleft{}one that indicates\textquotedblright{} (Monier-Williams).
\end{cousinweb}

\begin{grammarlens}[title={What you've built}]
One more word for getting around.
\end{grammarlens}

\subsection*{Your turn}
\begin{itemize}
\raggedright
  \item \emph{Say it:} \emph{sūcika}
  \item \emph{Say it:} \emph{sūcika}, once more
  \item \emph{Say it:} \emph{ā sūcika mīda nagarāniki dūraṁ rāsi undi.} — The distance to the city is written on that signpost.
  \item \emph{Recall:} say the Telugu for to break a rule, then the Telugu for a compass direction, then say \emph{sūcika} again
\end{itemize}

\begin{tcolorbox}[breakable,colback=teal!4,colframe=teal!35!black,title={Before you move on}]
What does \te{సూచిక} mean? (\textquotedblleft{}A signpost; an indicator\textquotedblright{}.) Say it once more.
\end{tcolorbox}
\section[\texorpdfstring{\te{దారి} \te{తప్పు} (dāri tappu)}{దారి తప్పు (dāri tappu)}]{\te{దారి} \te{తప్పు} (dāri tappu) — to lose one's way}

\label{lesson:TE-C448-dari-tappu}



\begin{quote}
Before the new one: say the Telugu for a compass direction, then the Telugu for a signpost.
\end{quote}

\subsection*{You'll want to know: \te{దారి} \te{తప్పు}}
\textbf{\te{దారి} \te{తప్పు}} — \emph{dāri tappu} — \textquotedblleft{}to lose one's way\textquotedblright{}.

\begin{quote}
\te{చీకటిలో} \te{మేము} \te{దారి} \te{తప్పాం}.
\end{quote}

\begin{quote}
\emph{cīkaṭilō mēmu dāri tappāṁ.}
\end{quote}

\begin{quote}
We lost our way in the dark.
\end{quote}

Past (I …): \te{దారి} \te{తప్పాను}, \emph{dāri tappānu}.

Word for word: to miss the way. \te{తప్పిపోవు}, to get lost, is close; \te{దారి} \te{తప్పు} is to take a wrong road.

\begin{cousinweb}
\textbf{\te{దారి} \te{తప్పు}} (\emph{dāri tappu}) is native Dravidian: \te{దారి}, \textquotedblleft{}a way, a road\textquotedblright{} (DEDR 3170), and \te{తప్పు}, \textquotedblleft{}to slip, to miss, to go astray\textquotedblright{} (DEDR 3071).
\end{cousinweb}

\begin{grammarlens}[title={What you've built}]
One more word for getting around.
\end{grammarlens}

\subsection*{Your turn}
\begin{itemize}
\raggedright
  \item \emph{Say it:} \emph{dāri tappu}
  \item \emph{Say it:} \emph{dāri tappu}, once more
  \item \emph{Say it:} \emph{cīkaṭilō mēmu dāri tappāṁ.} — We lost our way in the dark.
  \item \emph{Recall:} say the Telugu for a compass direction, then the Telugu for a signpost, then say \emph{dāri tappu} again
\end{itemize}

\begin{tcolorbox}[breakable,colback=teal!4,colframe=teal!35!black,title={Before you move on}]
What does \te{దారి} \te{తప్పు} mean? (\textquotedblleft{}To lose one's way\textquotedblright{}.) Say it once more.
\end{tcolorbox}
\section[\texorpdfstring{\te{మలుపు} \te{తిరుగు} (malupu tirugu)}{మలుపు తిరుగు (malupu tirugu)}]{\te{మలుపు} \te{తిరుగు} (malupu tirugu) — to take a turn}

\label{lesson:TE-C448-malupu-tirugu}



\begin{quote}
Before the new one: say the Telugu for a signpost, then the Telugu for to lose one's way.
\end{quote}

\subsection*{You'll want to know: \te{మలుపు} \te{తిరుగు}}
\textbf{\te{మలుపు} \te{తిరుగు}} — \emph{malupu tirugu} — \textquotedblleft{}to take a turn\textquotedblright{}.

\begin{quote}
\te{బస్సు} \te{ఎడమ} \te{వైపు} \te{మలుపు} \te{తిరిగింది}.
\end{quote}

\begin{quote}
\emph{bassu eḍama vaipu malupu tirigindi.}
\end{quote}

\begin{quote}
The bus turned to the left.
\end{quote}

Past (I …): \te{మలుపు} \te{తిరిగాను}, \emph{malupu tirigānu}.

\te{మలుపు}, a bend, + \te{తిరుగు}, to turn. It is also said of a story or a life: \te{కథ} \te{మలుపు} \te{తిరిగింది}, the story took a turn.

\begin{cousinweb}
\textbf{\te{మలుపు} \te{తిరుగు}} (\emph{malupu tirugu}) is a Telugu phrase: \te{మలుపు}, \textquotedblleft{}a bend, a turning\textquotedblright{}, probably native Dravidian (DEDR 4734 lists Telugu \emph{malapu}, \textquotedblleft{}to turn back\textquotedblright{}, and \emph{maluku}, \textquotedblleft{}a turn\textquotedblright{}, but not \emph{malupu} itself), and the native \te{తిరుగు}, \textquotedblleft{}to turn, to go round\textquotedblright{} (DEDR 3246).
\end{cousinweb}

\begin{grammarlens}[title={What you've built}]
One more word for getting around.
\end{grammarlens}

\subsection*{Your turn}
\begin{itemize}
\raggedright
  \item \emph{Say it:} \emph{malupu tirugu}
  \item \emph{Say it:} \emph{malupu tirugu}, once more
  \item \emph{Say it:} \emph{bassu eḍama vaipu malupu tirigindi.} — The bus turned to the left.
  \item \emph{Recall:} say the Telugu for a signpost, then the Telugu for to lose one's way, then say \emph{malupu tirugu} again
\end{itemize}

\begin{tcolorbox}[breakable,colback=teal!4,colframe=teal!35!black,title={Before you move on}]
What does \te{మలుపు} \te{తిరుగు} mean? (\textquotedblleft{}To take a turn\textquotedblright{}.) Say it once more.
\end{tcolorbox}
\section[\texorpdfstring{\te{అడ్డదారి} (aḍḍadāri)}{అడ్డదారి (aḍḍadāri)}]{\te{అడ్డదారి} (aḍḍadāri) — a shortcut}

\label{lesson:TE-C448-addadari}



\begin{quote}
Before the new one: say the Telugu for to lose one's way, then the Telugu for to take a turn.
\end{quote}

\subsection*{You'll want to know: \te{అడ్డదారి}}
\textbf{\te{అడ్డదారి}} — \emph{aḍḍadāri} — \textquotedblleft{}a shortcut\textquotedblright{}.

\begin{quote}
\te{అడ్డదారిలో} \te{వెళ్తే} \te{పది} \te{నిమిషాలు} \te{ముందు} \te{చేరతాం}.
\end{quote}

\begin{quote}
\emph{aḍḍadārilō veḷtē padi nimiṣālu mundu cēratāṁ.}
\end{quote}

\begin{quote}
If we take the shortcut we will get there ten minutes earlier.
\end{quote}

\te{అడ్డ}, across, + \te{దారి}, a way: a way that cuts across. It also means a dishonest shortcut.

\begin{cousinweb}
\textbf{\te{అడ్డదారి}} (\emph{aḍḍadāri}) is native Dravidian: \te{అడ్డ}, \textquotedblleft{}crosswise, across\textquotedblright{} (DEDR 83), and \te{దారి}, \textquotedblleft{}a way\textquotedblright{} (DEDR 3170).
\end{cousinweb}

\begin{grammarlens}[title={What you've built}]
One more word for getting around.
\end{grammarlens}

\subsection*{Your turn}
\begin{itemize}
\raggedright
  \item \emph{Say it:} \emph{aḍḍadāri}
  \item \emph{Say it:} \emph{aḍḍadāri}, once more
  \item \emph{Say it:} \emph{aḍḍadārilō veḷtē padi nimiṣālu mundu cēratāṁ.} — If we take the shortcut we will get there ten minutes earlier.
  \item \emph{Recall:} say the Telugu for to lose one's way, then the Telugu for to take a turn, then say \emph{aḍḍadāri} again
\end{itemize}

\begin{tcolorbox}[breakable,colback=teal!4,colframe=teal!35!black,title={Before you move on}]
What does \te{అడ్డదారి} mean? (\textquotedblleft{}A shortcut\textquotedblright{}.) Say it once more.
\end{tcolorbox}
